\chapter{Về Tương Lai}
\markboth{Về Tương Lai}{On the Future}

\section{Những gì em học hôm nay có thể giúp em sau 10 năm nữa.}
\begin{truyen}{Học hôm nay, dùng sau 10 năm}{What You Learn Today May Help in 10 Years}
\chuhoa{T}{hầy nói: ``Em không biết sau 10 năm em sẽ làm nghề gì, gặp ai, cần dùng kiến thức nào. Nhưng những gì em học hôm nay — cách tư duy, cách đọc, cách tính, cách viết — có thể xuất hiện lại khi em cần. Học không phải chỉ cho bài kiểm tra; học là gửi vốn cho tương lai. Sau 10 năm nữa em có thể cám ơn chính mình hôm nay.''}
\textit{(The teacher said: ``You don't know what you'll do, who you'll meet, or what knowledge you'll need in ten years. But what you learn today — how to think, read, calculate, write — may show up when you need it. Learning isn't just for the test; learning is investing in the future. In ten years you may thank today's you.'')}
\end{truyen}

\begin{baihoc}
Kiến thức và kỹ năng hôm nay là vốn; tương lai sẽ dùng đến nhiều hơn em tưởng.
\textit{(Today's knowledge and skills are capital; the future will use them more than you think.)}
\end{baihoc}

\begin{ghinhoanh}
\textbf{Short English to remember:}

What you learn today may help you in ten years.

\end{ghinhoanh}

\begin{tuhocnhanh}
1. Em có tin rằng học hôm nay có ích cho tương lai xa không? \textit{Do you believe today's learning helps your distant future?}

2. Kỹ năng nào em đang học có thể dùng lâu dài? \textit{Which skill you're learning could be used long term?}
\end{tuhocnhanh}
\ngancach

\section{Tương lai không được xây trong một ngày, mà trong rất nhiều ngày nhỏ.}
\begin{truyen}{Nhiều ngày nhỏ}{Many Small Days}
\chuhoa{T}{hầy nói: ``Em muốn tương lai giỏi, thành đạt — nhưng tương lai ấy không xây trong một ngày. Nó được xây bằng rất nhiều ngày nhỏ: hôm nay em học một tiết, mai em ôn một chút, ngày kia em làm thêm vài bài. Mỗi ngày nhỏ là một viên gạch. Em đặt đều tay, tương lai sẽ thành hình.''}
\textit{(The teacher said: ``You want a future where you're capable and successful — but that future isn't built in one day. It's built with many small days: today you study one lesson, tomorrow you review a bit, the next day you do a few more problems. Each small day is a brick. Lay them steadily and the future will take shape.'')}
\end{truyen}

\begin{baihoc}
Tương lai là tổng của rất nhiều ''hôm nay''; mỗi hôm nay đều đáng đầu tư.
\textit{(The future is the sum of many ``todays''; each today is worth investing in.)}
\end{baihoc}

\begin{ghinhoanh}
\textbf{Short English to remember:}

The future isn't built in one day, but in many small days.

\end{ghinhoanh}

\begin{tuhocnhanh}
1. ''Ngày nhỏ'' hôm nay em đang xây gì cho tương lai? \textit{What ``small day'' are you building for the future today?}

2. Em có đang bỏ qua ''ngày nhỏ'' vì nghĩ còn nhiều thời gian không? \textit{Do you skip ``small days'' thinking you have plenty of time?}
\end{tuhocnhanh}
\ngancach

\section{Thói quen hôm nay sẽ trở thành tính cách ngày mai.}
\begin{truyen}{Thói quen thành tính cách}{Habits Become Character}
\chuhoa{T}{hầy nhắc: ``Em hay trễ giờ, hay để dồn bài, hay nói dối nhỏ — đó là thói quen. Nhưng thói quen lặp lại sẽ thành tính cách. Người ta sẽ nói: em ấy là người hay trễ, hay để việc đến phút chót, không đáng tin. Thói quen hôm nay sẽ trở thành tính cách ngày mai. Em muốn tính cách thế nào, hãy rèn thói quen ấy từ bây giờ.''}
\textit{(The teacher reminded: ``You're often late, often leave work to the last minute, often tell small lies — those are habits. But habits repeated become character. People will say: that one is always late, always leaves things to the last minute, not reliable. Today's habits become tomorrow's character. If you want a certain character, build those habits now.'')}
\end{truyen}

\begin{baihoc}
Tính cách không tự nhiên có; nó được tạo từ thói quen lặp lại mỗi ngày.
\textit{(Character doesn't appear by itself; it's made from habits repeated every day.)}
\end{baihoc}

\begin{ghinhoanh}
\textbf{Short English to remember:}

Today's habits become tomorrow's character.

\end{ghinhoanh}

\begin{tuhocnhanh}
1. Thói quen nào của em đang hình thành ''tính cách''? \textit{Which habits of yours are forming ``character''?}

2. Em muốn tính cách mình thế nào — và cần thói quen gì? \textit{What character do you want — and what habits do you need?}
\end{tuhocnhanh}
\ngancach

\section{Người đọc nhiều hôm nay sẽ nói hay ngày mai.}
\begin{truyen}{Đọc nhiều, nói hay}{Read More Today, Speak Well Tomorrow}
\chuhoa{T}{hầy nói với em hay ngại phát biểu: ``Em muốn nói hay, nói đúng ý — thì hôm nay em phải đọc nhiều, nghe nhiều, suy nghĩ nhiều. Người đọc nhiều có vốn từ và ý tưởng; khi cần nói hoặc viết, họ có cái để lấy ra. Người đọc nhiều hôm nay sẽ nói hay ngày mai. Em đừng chờ đến khi cần mới đọc.''}
\textit{(The teacher told a student who was shy to speak: ``If you want to speak well and clearly — you need to read, listen, and think a lot today. Those who read a lot have a store of words and ideas; when they need to speak or write, they have something to draw on. Those who read a lot today will speak well tomorrow. Don't wait until you need it to start reading.'')}
\end{truyen}

\begin{baihoc}
Đọc là nạp vào; nói và viết là lấy ra — không nạp thì không có gì để lấy.
\textit{(Reading is input; speaking and writing are output — no input, nothing to give out.)}
\end{baihoc}

\begin{ghinhoanh}
\textbf{Short English to remember:}

Those who read a lot today will speak well tomorrow.

\end{ghinhoanh}

\begin{tuhocnhanh}
1. Em có đọc đều — sách, bài — mỗi ngày không? \textit{Do you read regularly — books, articles — every day?}

2. Em có thấy đọc giúp em nói và viết rõ hơn không? \textit{Do you notice reading helping you speak and write more clearly?}
\end{tuhocnhanh}
\ngancach

\section{Nếu em muốn đi xa, hãy chuẩn bị từ bây giờ.}
\begin{truyen}{Chuẩn bị từ bây giờ}{Prepare Now}
\chuhoa{T}{hầy nói: ``Em muốn đi xa — xa về tri thức, xa về nghề nghiệp, xa về cuộc sống. Không ai đi xa được nếu đứng yên đến phút chót rồi mới chạy. Nếu em muốn đi xa, hãy chuẩn bị từ bây giờ: học đều, rèn thói quen, giữ sức khỏe, giữ quan hệ tốt. Chuẩn bị hôm nay là bước đầu của con đường xa.''}
\textit{(The teacher said: ``You want to go far — far in knowledge, career, life. No one goes far by standing still until the last minute and then running. If you want to go far, prepare now: study steadily, build habits, take care of your health, keep good relationships. Today's preparation is the first step on the long road.'')}
\end{truyen}

\begin{baihoc}
Đi xa không bắt đầu khi em ''lên đường''; nó bắt đầu từ những chuẩn bị hôm nay.
\textit{(Going far doesn't start when you ``hit the road''; it starts with today's preparation.)}
\end{baihoc}

\begin{ghinhoanh}
\textbf{Short English to remember:}

If you want to go far, prepare from now.

\end{ghinhoanh}

\begin{tuhocnhanh}
1. Em đang chuẩn bị gì ''từ bây giờ'' cho tương lai? \textit{What are you preparing ``from now'' for the future?}

2. Điều gì sẽ giúp em ''đi xa'' — và em đã bắt đầu chưa? \textit{What will help you ``go far'' — and have you started?}
\end{tuhocnhanh}
\ngancach

\section{Người thành công thường là người bắt đầu sớm hơn người khác.}
\begin{truyen}{Bắt đầu sớm}{Starting Early}
\chuhoa{T}{hầy kể: ``Hai người cùng muốn học ngoại ngữ. Một người bắt đầu từ năm lớp 10, mỗi ngày 15 phút. Người kia đợi đến khi cần mới học. Đến lúc cần — đi du học, xin việc — người bắt đầu sớm đã có vốn sẵn; người bắt đầu muộn phải chạy gấp. Thành công thường thuộc về người bắt đầu sớm hơn.''}
\textit{(The teacher said: ``Two people both wanted to learn a language. One started in grade 10, fifteen minutes a day. The other waited until they needed it. When the need came — study abroad, job application — the early starter had the capital ready; the late one had to rush. Success often goes to those who start earlier.'')}
\end{truyen}

\begin{baihoc}
Bắt đầu sớm cho em thời gian tích lũy; bắt đầu muộn chỉ cho em áp lực.
\textit{(Starting early gives you time to build; starting late only gives you pressure.)}
\end{baihoc}

\begin{ghinhoanh}
\textbf{Short English to remember:}

Successful people often simply start earlier than others.

\end{ghinhoanh}

\begin{tuhocnhanh}
1. Việc gì em đang ''đợi'' mới làm — có nên bắt đầu sớm không? \textit{What are you ``waiting'' to do — should you start early?}

2. Em có thể ''bắt đầu sớm'' điều gì ngay từ học kỳ này? \textit{What can you ``start early'' this semester?}
\end{tuhocnhanh}
\ngancach

\section{Mỗi cuốn sách em đọc là một cánh cửa mới.}
\begin{truyen}{Mỗi cuốn sách là một cánh cửa}{Every Book Is a New Door}
\chuhoa{T}{hầy nói: ``Mỗi cuốn sách em đọc — dù truyện, dù khoa học, dù lịch sử — mở cho em một thế giới hoặc một cách nghĩ mới. Em không cần đọc hàng trăm cuốn; em chỉ cần đọc đều và đọc có suy nghĩ. Mỗi cuốn sách là một cánh cửa. Em mở càng nhiều, em càng thấy đường đi rộng hơn.''}
\textit{(The teacher said: ``Every book you read — whether story, science, or history — opens a new world or a new way of thinking for you. You don't need to read hundreds; you need to read steadily and think. Every book is a door. The more you open, the wider the path you see.'')}
\end{truyen}

\begin{baihoc}
Sách mở thêm cánh cửa; không đọc là chỉ đi trong vài căn phòng quen.
\textit{(Books open more doors; not reading is staying in a few familiar rooms.)}
\end{baihoc}

\begin{ghinhoanh}
\textbf{Short English to remember:}

Every book you read is a new door.

\end{ghinhoanh}

\begin{tuhocnhanh}
1. Cuốn sách gần đây đã ''mở'' cho em điều gì? \textit{What did a recent book ``open'' for you?}

2. Em có đặt mục tiêu đọc sách đều không? \textit{Do you set a goal to read regularly?}
\end{tuhocnhanh}
\ngancach

\section{Đừng hỏi tương lai cho em điều gì, hãy hỏi em đang chuẩn bị gì cho tương lai.}
\begin{truyen}{Chuẩn bị gì cho tương lai}{What Are You Preparing for the Future}
\chuhoa{T}{hầy nói: ``Nhiều em hỏi: tương lai em sẽ làm gì, sẽ giàu không, sẽ hạnh phúc không. Câu hỏi đúng hơn là: em đang chuẩn bị gì cho tương lai? Em đang học, đang rèn thói quen, đang giữ sức khỏe và quan hệ không? Tương lai không trả lời khi em chỉ hỏi; tương lai trả lời khi em hành động hôm nay.''}
\textit{(The teacher said: ``Many ask: what will I do in the future, will I be rich, will I be happy. The better question is: what are you preparing for the future? Are you studying, building habits, taking care of health and relationships? The future doesn't answer when you only ask; it answers when you act today.'')}
\end{truyen}

\begin{baihoc}
Tương lai không cho em gì nếu em không cho tương lai gì — bằng cách chuẩn bị hôm nay.
\textit{(The future gives you nothing if you give the future nothing — by preparing today.)}
\end{baihoc}

\begin{ghinhoanh}
\textbf{Short English to remember:}

Don't ask what the future will give you; ask what you're preparing for the future.

\end{ghinhoanh}

\begin{tuhocnhanh}
1. Em thường hỏi ''tương lai sẽ thế nào'' hay ''mình đang chuẩn bị gì''? \textit{Do you usually ask ``what will the future be'' or ``what am I preparing''?}

2. Em đang ''cho'' tương lai điều gì bằng hành động hôm nay? \textit{What are you ``giving'' the future by today's actions?}
\end{tuhocnhanh}
\ngancach

\section{Thời gian của tuổi trẻ là vốn quý nhất.}
\begin{truyen}{Vốn quý nhất}{The Most Precious Capital}
\chuhoa{T}{hầy nói: ``Tiền có thể kiếm lại, đồ có thể mua lại. Nhưng thời gian tuổi trẻ — năm tháng em đang ở trường, đầu óc còn nhanh, cơ thể còn khỏe — không quay lại. Mỗi ngày em bỏ qua là mất đi vĩnh viễn. Thời gian tuổi trẻ là vốn quý nhất. Em dùng nó thế nào sẽ quyết định em có gì khi lớn.''}
\textit{(The teacher said: ``Money can be earned again, things can be bought again. But the time of youth — the years you're in school, mind still sharp, body still strong — doesn't come back. Every day you waste is lost forever. The time of youth is the most precious capital. How you use it will decide what you have when you grow up.'')}
\end{truyen}

\begin{baihoc}
Tuổi trẻ là vốn không tái tạo; tiêu vào học và rèn là đầu tư đúng.
\textit{(Youth is non-renewable capital; spending it on learning and practice is the right investment.)}
\end{baihoc}

\begin{ghinhoanh}
\textbf{Short English to remember:}

The time of youth is your most precious capital.

\end{ghinhoanh}

\begin{tuhocnhanh}
1. Em có đang ''tiêu'' thời gian tuổi trẻ vào việc có ích không? \textit{Are you ``spending'' your youth on what's useful?}

2. Một ngày ''mất'' với em nghĩa là gì? \textit{What does a ``wasted'' day mean to you?}
\end{tuhocnhanh}
\ngancach

\section{Một ngày nào đó em sẽ hiểu: việc học hôm nay là món quà lớn nhất cho chính mình.}
\begin{truyen}{Món quà lớn nhất cho chính mình}{The Greatest Gift to Yourself}
\chuhoa{T}{hầy nói cuối năm: ``Bây giờ có thể em thấy học là gánh nặng, là áp lực. Nhưng một ngày nào đó — khi em cần dùng kiến thức, khi em đứng trước cơ hội, khi em nhìn lại — em sẽ hiểu: việc học hôm nay là món quà lớn nhất em dành cho chính mình. Không ai cho em món quà ấy; chỉ có em có thể tự cho mình bằng cách học.''}
\textit{(The teacher said at year's end: ``Right now you may see school as a burden, as pressure. But one day — when you need that knowledge, when you stand before an opportunity, when you look back — you'll understand: learning today is the greatest gift you give yourself. No one else can give you that gift; only you can by learning.'')}
\end{truyen}

\begin{baihoc}
Học là món quà mình tự cho mình; không ai thay em mở món quà ấy được.
\textit{(Learning is the gift you give yourself; no one else can open that gift for you.)}
\end{baihoc}

\begin{ghinhoanh}
\textbf{Short English to remember:}

One day you'll understand: learning today is the greatest gift you give yourself.

\end{ghinhoanh}

\begin{tuhocnhanh}
1. Em có tin rằng học là ''món quà cho chính mình'' không? \textit{Do you believe learning is ``a gift to yourself''?}

2. Em muốn ''món quà'' mình tự cho mình là gì — kiến thức, thói quen, tính cách? \textit{What ``gift'' do you want to give yourself — knowledge, habits, character?}
\end{tuhocnhanh}
\ngancach
